\documentclass[10pt]{article}  

%%%%%%%% Preamble %%%%%%%%%%%%
\title{Degree project}
\usepackage[utf8]{inputenc} % File coding uses utf8
\usepackage{amsmath} % Extra commands for math
\usepackage{amssymb} % Math symbols 
\usepackage{graphicx} % Include images in LaTeX
\usepackage{color} % Coloring text
\usepackage{subfigure} % Manage multiple figures
\usepackage{float} % Allow you to use [H] specifier to force the position of the images 
\usepackage{capt-of} % Defines a command \captionof for putting a caption to something that’s not a float.
\usepackage{sidecap} % Defines environments called SCfigure and SCtable (analogous to figure and table) to typeset captions sideways
	\sidecaptionvpos{figure}{c} % Alignment
\usepackage{caption} % to customize the captions in floating environments like figure and table
\usepackage{commath} % Mathematics typesetting support
\usepackage{cancel} % Place lines through maths formulae
\usepackage{anysize} % to set up document margins
\marginsize{2cm}{2cm}{2cm}{2cm} % Left, right, up, down
\usepackage{appendix} %Extra control of appendices

% Refereces as links with colors
\usepackage[colorlinks=true,plainpages=true,citecolor=blue,linkcolor=blue]{hyperref}

% Header and Footer
\usepackage{fancyhdr} 
\pagestyle{fancy}
\fancyhf{}
\fancyhead[L]{\footnotesize Faculty of Computer and Information Technology} 
\fancyhead[R]{\footnotesize Jordan University of Science and Technology}   
\fancyfoot[R]{\footnotesize Graduation Project 1}  
\fancyfoot[C]{\thepage}  % center
\fancyfoot[L]{\footnotesize Network Engineering and Security Department}  %izquierda
\renewcommand{\footrulewidth}{0.4pt}
\fancypagestyle{firststyle}
{
   \fancyhf{}
}

\usepackage{listings} % To use source code
\definecolor{dkgreen}{rgb}{0,0.6,0} % Color for using code
\definecolor{gray}{rgb}{0.5,0.5,0.5} 
% Language to use
\lstset{language=Matlab,
   keywords={break,case,catch,continue,else,elseif,end,for,function,
      global,if,otherwise,persistent,return,switch,try,while},
   basicstyle=\ttfamily,
   keywordstyle=\color{blue},
   commentstyle=\color{red},
   stringstyle=\color{dkgreen},
   numbers=left,
   numberstyle=\tiny\color{gray},
   stepnumber=1,
   numbersep=10pt,
   backgroundcolor=\color{white},
   tabsize=4,
   showspaces=false,
   showstringspaces=false}

\title{Degree project}

%%%%%%%% Preamble ends %%%%%%%%%%%%

\begin{document}

%%%%%%%%%%%%%%%%%%%%%%%%%%%%%%%% Cover Page %%%%%%%%%%%%%%%%%%%%%%%%%%%%%%%%%%%%%%%%%%%

\thispagestyle{firststyle}
\begin{center}
  \begin{minipage}{0.48\textwidth} \begin{flushleft}
    \includegraphics[scale = 0.75]{images/Just_logo.png}
  \end{flushleft}\end{minipage}
  \begin{minipage}{0.48\textwidth} \begin{flushright}
    \includegraphics[scale = 0.75]{images/Just_logo.png}
  \end{flushright}\end{minipage}

  \vspace*{-2.5cm}
  \textsc{\huge Jordan University \\
  \vspace{5px} of Science \\
  \vspace{15px} and Technology}\\

  \vspace*{1.5cm}

  \textsc{\LARGE Faculty of Computer and Information Technology}\\[1.5cm]
  \textsc{\LARGE Network Engineering and Security Department}\\[1.5cm]

  \begin{minipage}{0.9\textwidth} 
    \begin{center}
      \textsc{\LARGE Graduation Project 1 Report}
    \end{center}
  \end{minipage}\\[0.5cm]
  
  \vspace*{1cm}

  { \huge \bfseries Title}\\[0.4cm]	
  \vspace*{1cm}
  { \large 
    \emph{Authors:} \\	
      Author 1, University ID \\
      Author 2, University ID \\
      Author 3, University ID \\
      Author 4, University ID \\
    \vspace*{0.5cm}
    \emph{Advisers:} \\													  
      Advisor 1 \\
      Advisor 2	\\
      Advisor 3	\\
      Advisor 4	\\
  }


  \vspace{2cm} 	
  \begin{center}
    {Urcuquí - \today}
  \end{center}
  
\end{center}
																		
\newpage																		
%%%%%%%%%%%%%%%%%%%% Cover page ends %%%%%%%%%%%%%%%%%%%%%%%%%%%%%%%%

\tableofcontents 

\newpage

\section{Introduction}

In this section you are required to present the idea behind your project. A general description of the problem you are solving must be provided. You need to discuss the challenges you will be facing. Also, you need to describe the goals of the project (what you are trying to achieve).  The targeted users for your project should be discussed (i.e. the project is oriented toward helping the elderly, college students, terminal patients, etc.).


 
\section{Background}
In this section you are required to provide a brief history of the work you are planning to do. You should provide a description of the algorithms, tools (hardware and software), and development environment you are using to realize your project goals. For example, if your work is cryptography oriented, you need to provide a background for the ciphers you are using or modifying. Note that all external resources must be cited similar to the one found in the references section \cite{IEEEreferencias:Ref1}. and \cite{1989:Rabiner}.

\section{System Model}
Start this section by detailing your problem statement (i.e.; what problem you intend to address and why). Then, explain the proposed system model in details. You are expected to present the steps to be followed in building your system in order to achieve its goals. Your model can be a diagram, a flowchart, or a pseudocode for the procedure you are using. Include the model assumptions, requirements, and constraints. 

\subsection{General Objective}
bla bla 

\subsection{Specific Objectives}
bla bla 

\section{Completed Tasks}

[For this section, provide details about the tasks you have completed thus far in your project. Mention the major milestones you have reached and also discuss the limitations that you have faced. Examples of items to be included are building a major project component, successfully running an attack, or implementing a compression algorithm. On the other hand, installing a compiler or downloading a tool are not considered milestones and should not be discussed. \\

Along with these points, it is necessary to discuss how these steps would lead to the successful completion of the project or how some limitations would have an adverse effect and the methods you have applied or intend to apply to overcome them. \\

Please adhere to the 3-page limit (excluding the cover page).


%%%%%%% Bibliografía %%%%%%%%
\bibliographystyle{bst/IEEEtran} 
\addcontentsline{toc}{section}{References}  
\bibliography{bib/IEEEreferences} 
%%%%%%% Bibliografía %%%%%%%%    

\end{document}